\documentclass[11pt]{article}
\usepackage[T1]{fontenc}
\usepackage[utf8]{inputenc}
\usepackage{lmodern,amsmath,amssymb,amsthm,mathtools,booktabs,microtype}
\usepackage[margin=1in]{geometry}
\usepackage[hidelinks]{hyperref}
\setlength{\emergencystretch}{3em}
\newtheorem{theorem}{Theorem}[section]
\newtheorem{lemma}[theorem]{Lemma}
\newtheorem{proposition}[theorem]{Proposition}
\newtheorem{corollary}[theorem]{Corollary}
\theoremstyle{remark}\newtheorem{remark}[theorem]{Remark}
\newcommand{\R}{\mathbb R}\newcommand{\Q}{\mathbb Q}\newcommand{\Z}{\mathbb Z}
\newcommand{\CP}{\mathcal{CP}}\newcommand{\COP}{\mathcal{COP}}
\newcommand{\Ry}{\mathcal R}\newcommand{\tr}{\operatorname{tr}}
\newcommand{\diag}{\operatorname{diag}}\newcommand{\cone}{\operatorname{cone}}
\newcommand{\ip}[2]{\langle #1,#2\rangle}\newcommand{\trans}{^{\mathsf T}}
\newcommand{\norm}[1]{\lVert #1\rVert}
\title{Exponential local pivot complexity for ordinary completely positive separation\\[3pt]\large Every descending adjacent rule, canonical initialization, and robust positive inputs}
\author{Manuscript II research continuation}
\date{15 November 2026}
\begin{document}\maketitle
\begin{abstract}
For each fixed order $n\ge5$ we construct strictly entrywise-positive, positive-definite rational matrices outside the ordinary completely positive cone on which every strictly descending adjacent-vertex rule, initialized at the canonical tridiagonal perfect form, requires exponentially many pivots in the binary input-and-distance parameter. The lower bound holds on an explicit Frobenius ball and does not require uniqueness of the improving ray. A dimension-independent bound of $5/2$ on the movement of each cross-column coordinate along a bounded edge of a fixed-positive-definite-block Ryshkov face replaces trajectory-specific analysis. An arithmetic gap and a cofactor height bound prevent every improving extreme ray from leaving that face during the entire prefix. No scalar identity correction of any improving ray separates the inputs there. In contrast, a strictly copositive rational separator with logarithmic coordinate height is explicit. The proof is for ordinary $\CP_n$, without additional column restrictions, and uses the standard canonical initialization without preprocessing. The constants are conservative; historical first priority and external proof certification are not asserted.
\end{abstract}

\section{Model and main statement}
The pairing is $\ip A P=\tr(AP)$. Set
\[
 \CP_n=\cone\{xx\trans:x\in\R_+^n\},\qquad
 \COP_n=\{P:P[x]\ge0\text{ for }x\ge0\},
\]
where $P[x]=x\trans Px$, and
\[
 \Ry_n=\{P\in\mathbb S^n:P[v]\ge1\text{ for all }0\ne v\in\Z_+^n\}.
\]
Let $T_j$ have diagonal one, adjacent entries $-1/2$, and all other entries zero. A local rule in this paper starts at $T_n$, chooses any extreme tangent ray $R$ with $\ip A R<0$, and moves to the adjacent vertex if the ray is bounded. The standard certificate exits are $\ip A P<0$ at the current vertex or $R\in\COP_n$. We additionally allow testing $R+\tau I$ for arbitrary real $\tau$. No computational restriction is imposed on choosing among the improving rays; in particular the choice may be randomized, history dependent, or informed by the entire input.

This is the adjacent-pivot model of Procedure 3 in \cite{DSV}, augmented by the stated identity correction. An unrelated separator, an input-dependent starting vertex, nonadjacent jumps, or a non-descending move is outside the model. The existence of a short global separator is compatible with the lower bound precisely because it is a nonlocal bypass.

For $n\ge5$, $m\ge1$, define
\begin{equation}\label{eq:params}
 \begin{split}
 r&=n-1,\qquad d=n(n+1)/2,\qquad q=2n,\qquad Y=2q^r,\\
 U&=2Yn^2(m+1),\qquad E=(d-1)![2U^2]^{d-1},\\
 N&=20000n^3E^2,\qquad \epsilon=\frac1{100n^3},\qquad
 \rho=\frac1{1600n^4N^2}.
 \end{split}
\end{equation}
Let $G$ be the adjacency matrix of the five-cycle on the ordered coordinates $(1,2,3,4,n)$, zero elsewhere, and put
\begin{equation}\label{eq:targets}
 \begin{split}
 F&=I_n+\frac35G,\qquad D=\diag(1,8/5,8/5,1,\ldots,1),\\
 C_{ij}&=\min(i,j)(n+1-\max(i,j)),\\
 B&=DFD+\epsilon C,\qquad S_N=\diag(N,\ldots,N,1),\\
 B_N&=S_NBS_N,\qquad A_N=\frac{B_N}{4nN^2}.
 \end{split}
\end{equation}
All objects in these formulas are explicit and rational.

\begin{theorem}[Canonical, ordinary-cone, all-rule lower bound]\label{thm:main}
For every $n\ge5$, $m\ge1$, and every symmetric $A$ satisfying
\begin{equation}\label{eq:ball}\norm{A-A_N}_F\le\rho,
\end{equation}
the following statements hold.
\begin{enumerate}
\item $\norm A_F\le1$, $A\succ0$, every entry of $A$ is positive, and $A\notin\CP_n$.
\item Every local rule in the model above makes at least $m$ adjacent-vertex pivots from $T_n$. During these pivots an improving extreme ray always exists, every such ray is bounded, and neither standard certificate exit applies.
\item At every departure vertex in this prefix and for every improving extreme ray $R$, no real $\tau$ makes $R+\tau I$ a copositive separator of $A$.
\item One explicit rational $W_N$ separates the entire ball and satisfies
\[
 W_N[x]\ge\frac{\norm x_2^2}{20nN^2}\quad(x\ge0),\qquad
 \norm{W_N}_F\le6,\qquad
 \ip{W_N}A\le-\frac1{6nN^2}.
\]
Consequently $\operatorname{dist}_F(A,\CP_n)\ge1/(36nN^2)$.
\end{enumerate}
For each fixed $n$, the rational centers have total input length $O_n(\log(m+1))$ and $W_N$ has coordinate bit height $O_n(\log(m+1))$. If $L_m$ is the center's total rational encoding length and
\[
 H_m=L_m+\left\lceil\log_2\frac1{\operatorname{dist}_F(A_N,\CP_n)}\right\rceil+2,
\]
then $H_m=\Theta_n(\log(m+1))$. Thus the required number of pivots is $2^{\Omega_n(H_m)}$.
\end{theorem}

The bit-length assertion is for the explicit centers. Rational perturbations may have arbitrarily long encodings, even inside a small ball. The theorem does not assert that every rule eventually terminates; it asserts an unavoidable finite prefix before any of the allowed exits. It also does not assert a unique path. All possible improving choices are covered.

\section{A bounded-movement theorem for fixed positive-definite blocks}
This section supplies the geometric ingredient independently of the constructed objective.

Let $H\in\mathbb S^r$ be positive definite. Consider the epigraph defined by
\begin{equation}\label{eq:g}
 g_H(b)=\sup_{u\in\Z_+^r,\ y\in\Z_{\ge1}}
 \frac{1-u\trans Hu-2y b\trans u}{y^2},\qquad
 \mathcal E_H=\{(b,c):c\ge g_H(b)\}.
\end{equation}
Fix any integer $q\ge2$ with $\norm H_2r/q^2\le1/2$ and let
\[
 \kappa=\frac1{2q^{2r}},\qquad Y=2q^r.
\]
For each $b$, let $z$ be the unique minimizer of $x\trans Hx+2b\trans x$ on $\R_+^r$ and set $\ell=b+Hz$. The optimality conditions give
\begin{equation}\label{eq:kkt}
 z\ge0,\quad\ell\ge0,\quad z_i\ell_i=0.
\end{equation}

\begin{lemma}[Uniform transverse gap and active-height control]\label{lem:gap}
For every $b$,
\begin{equation}\label{eq:gap}
 z\trans Hz+\kappa\le g_H(b)\le z\trans Hz+1.
\end{equation}
The supremum is attained. At an active constraint $(u,y)$ on the graph of $g_H$, with $e=u-yz$,
\begin{equation}\label{eq:active}
 y\le Y,\qquad e\trans He+2y\ell\trans u
 =1-y^2\bigl(g_H(b)-z\trans Hz\bigr)\le1.
\end{equation}
The epigraph is locally polyhedral. At each of its vertices, $0\le\ell_i\le1/2$ for every $i$.
\end{lemma}
\begin{proof}
Completing the square gives
\begin{equation}\label{eq:complete}
 \frac{1-u\trans Hu-2yb\trans u}{y^2}
 =z\trans Hz+\frac{1-e\trans He-2y\ell\trans u}{y^2}.
\end{equation}
Its last numerator is at most one, proving the upper bound. On the positive support of $z$, simultaneous Dirichlet approximation supplies $1\le y\le q^r$ and integers $u_i$ with $|u_i-yz_i|\le1/q$. Set $u_i=0$ on the other coordinates. Each approximating integer is nonnegative: a negative integer would have error at least one. Moreover $\ell\trans u=0$. Hence $e\trans He\le\norm H_2r/q^2\le1/2$, proving the lower bound. If $z=0$, take $u=0,y=1$.

For completeness, the approximation follows by placing the $q^s+1$ points $\{kz_S\}$, $0\le k\le q^s$, in $q^s$ boxes of side $1/q$, where $s$ is the positive-support size. Subtract two points in one box.

In \eqref{eq:complete}, all denominators $y>Y$ give values strictly below the lower bound in \eqref{eq:gap}. For the remaining finitely many denominators, a value at least $z\trans Hz+\kappa$ requires $u-yz$ in a fixed ellipsoid. Thus only finitely many numerators matter; the supremum is attained, and \eqref{eq:active} follows. The same bounds on a neighborhood of $b$ prove local polyhedrality. Continuity of $z(b)$ follows, for example, from strong convexity of the minimized quadratic.

At a vertex $c=g_H(b)$ and the active affine normals $(2u/y,1)$ span $\R^{r+1}$. For each $i$, some active $u_i$ is therefore positive. Equation \eqref{eq:active} implies $2y\ell_i u_i\le1$, and hence $\ell_i\le1/2$.
\end{proof}

\begin{theorem}[Bounded coordinate movement along every edge]\label{thm:movement}
If $(b,c)$ and $(b',c')$ are endpoints of a bounded edge of $\mathcal E_H$, then
\begin{equation}\label{eq:movement}
 |b'_i-b_i|\le\frac12+2\sqrt{H_{ii}}\qquad(1\le i\le r).
\end{equation}
For $H=T_r$ this is $\norm{b'-b}_\infty\le5/2$, independently of dimension, location, degeneracy, and edge choice.
\end{theorem}
\begin{proof}
A bounded edge in a full-dimensional $(r+1)$-dimensional locally polyhedral epigraph has a common active constraint $(u,y)$ at its two endpoints (indeed at least $r$ independent common normals). At either endpoint,
\[
 b+H(u/y)=\ell+He/y.
\]
Cauchy--Schwarz in the $H$-inner product and \eqref{eq:active} give $|(He)_i|\le\sqrt{H_{ii}}$. By Lemma~\ref{lem:gap}, the $i$th coordinate lies in
\[
 [-\sqrt{H_{ii}}/y,\ 1/2+\sqrt{H_{ii}}/y].
\]
The same $u,y$ occur at both endpoints. Taking the length of this interval and using $y\ge1$ proves the assertion.
\end{proof}

\begin{lemma}[Recession directions]\label{lem:recession}
The recession cone of $\mathcal E_H$ is $\{(h,a):h\ge0,\ a\ge0\}$.
\end{lemma}
\begin{proof}
A recession direction must satisfy $a+2h\trans u/y\ge0$ for all $u\in\Z_+^r,y\ge1$. Taking $u=0$ gives $a\ge0$, and taking $u=ke_i,y=1$ with arbitrarily large $k$ gives $h_i\ge0$. These conditions are also sufficient.
\end{proof}

\section{The canonical face and arithmetic ray heights}
The identity
\begin{equation}\label{eq:root}
 T_j[x]=\frac12\left(x_1^2+x_j^2+\sum_{i=1}^{j-1}(x_i-x_{i+1})^2\right)
\end{equation}
shows that its nonnegative integer minimum is one, attained precisely at the consecutive-interval indicators. Their rank-one matrices are linearly independent: the diagonal entries are recovered from singleton intervals and the off-diagonals from four-term interval differences. Thus $T_j$ is a perfect vertex, with exactly $j(j+1)/2$ active normals.

Retain, at $T_n$, all interval equalities on the first $r=n-1$ coordinates. Their intersection with $\Ry_n$ is the face
\begin{equation}\label{eq:face}
 \mathcal F=\left\{P(b,c)=\begin{pmatrix}T_r&b\\b\trans&c\end{pmatrix}: c\ge g_{T_r}(b)\right\}.
\end{equation}
It has relative dimension $r+1$ and contains $T_n=P(-e_r/2,1)$. For $T_r$, $\norm{T_r}_2<2$, so the choices $q=2n$, $Y=2q^r$ in \eqref{eq:params} are valid in Lemma~\ref{lem:gap}.

Every point of this face is strictly copositive. Indeed, with $\ell,z$ as above and $a=c-z\trans T_rz\ge\kappa$,
\[
 P[(x,y)]=(x-yz)\trans T_r(x-yz)+2y\ell\trans x+ay^2>0
\]
for every nonzero $(x,y)\ge0$. Thus the full Ryshkov set is locally polyhedral at every point of $\mathcal F$: a sufficiently small neighborhood has a uniform positive copositivity margin, so only finitely many integer vectors can define constraints there. Standard local polyhedral edge and tangent-cone facts are consequently applicable. This also proves that a nonrecession extreme tangent ray in $\mathcal F$ reaches a genuine adjacent endpoint, rather than a nonattained boundary infimum.

\begin{lemma}[Heights in the first $m$ possible steps]\label{lem:heights}
Any vertex of $\mathcal F$ reachable from $T_n$ in at most $m$ bounded edges of $\mathcal F$ satisfies
\[
 \norm b_\infty\le3(m+1),\quad
 \norm z_\infty\le n^2(m+1),\quad
 0<c\le3n^5(m+1)^2,\quad
 \norm P_F\le10n^5(m+1)^2.
\]
Every active vector of the full $\Ry_n$ at this vertex has coordinates at most $U$ from \eqref{eq:params}. Every extreme tangent ray has a primitive integral generator $R$ with
\begin{equation}\label{eq:rayheight}
 |R_{ij}|\le E,\qquad \norm R_F\le nE.
\end{equation}
\end{lemma}
\begin{proof}
The movement theorem gives $\norm b_\infty\le1/2+5m/2\le3(m+1)$. The inverse formula
\[
 (T_r^{-1})_{ij}=\frac{2\min(i,j)(n-\max(i,j))}{n}
\]
gives row sums $i(n-i)\le n^2/4$ and diagonal entries at most $n/2$. Since $z=T_r^{-1}(\ell-b)$ and $0\le\ell_i\le1/2$, the displayed bound on $z$ follows. The upper bound on $g$ and $\norm{T_r}_2<2$ give the bounds on $c$ and $P$.

An active vector with last coordinate zero is a top interval indicator. Otherwise write it $(u,y)$. Lemma~\ref{lem:gap} gives $y\le Y$ and
\[
 |u_i-yz_i|\le\sqrt{(T_r^{-1})_{ii}}\le\sqrt{n/2}.
\]
Therefore $u_i\le Yn^2(m+1)+\sqrt{n/2}\le U$.

In upper-triangular coordinates an active equation has integer coefficients $v_i^2$ on the diagonal and $2v_iv_j$ off the diagonal, all of magnitude at most $2U^2$. At a vertex the tangent cone is full dimensional and pointed. Its extreme ray annihilates $d-1$ independent active equations. The cofactor formula produces an integral generator bounded by $(d-1)![2U^2]^{d-1}=E$. Dividing by a greatest common divisor can only reduce this bound. No simplicity assumption is used.
\end{proof}

\section{An ordinary completely positive obstruction compatible with the face}
Write the fixed matrix $B$ in \eqref{eq:targets} as
\[
 B=\begin{pmatrix}B_0&c_0\\c_0\trans&a_0\end{pmatrix}.
\]
\begin{lemma}[Fixed objective data]\label{lem:data}
The matrix $B$ is strictly entrywise positive, $B\succ\frac1{40}I$, and $\norm B_F\le3n$. Every entry is at least $\epsilon$. Also
\[
 0<a_0<2,\qquad c_0>0,\qquad \sum_i(c_0)_i<2.
\]
The leading block $B_0$ is a strictly positive combination of all rank-one top-interval normals, with every coefficient at least $\epsilon$. All entries of $B$ have common denominator $500n^3$, and $\ip{B_0}{T_r}\ge1$.
There is a rational $W_0\in\COP_n$ with $\norm{W_0}_F\le5$ and $\ip{W_0}B\le-19/20$.
\end{lemma}
\begin{proof}
The smallest eigenvalue of the five-cycle block of $F$ is $(7-3\sqrt5)/10>1/40$; the remaining eigenvalues outside the block are one. Since the diagonal entries of $D$ are at least one, $DFD\succ I/40$. The anchor $C$ is the sum of the squares of all interval indicators. It is positive definite, every entry is at least one, and $\norm C_F\le\tr C=n(n+1)(n+2)/6\le n^3$.

The leading block of $DFD$ has the following interval decomposition. The singleton weights on indices $1,2,3,4$ are respectively
\[
 1/25,\quad8/125,\quad8/125,\quad1/25,
\]
the weights on $[1,2],[2,3],[3,4]$ are $24/25,192/125,24/25$, and any other top coordinate has singleton weight one. The leading block of $C$ gives weight one to each top interval, plus one more to those ending at $r$. This proves the exposure assertion. The base weights alone sum to $r-42/125>1$.

Here $a_0=1+\epsilon n$, while $c_0$ has base values $3/5$ at indices $1,4$ and zero elsewhere, plus $\epsilon i$ at index $i$. Thus its sum is $6/5+\epsilon n(n-1)/2<2$. Every entry of $B$ is below three, proving the norm bound. The denominator claim follows because $125$ and $100n^3$ both divide $500n^3$.

Let $H_5=J_5-2\operatorname{Adj}(C_5)$ be the Horn matrix. It is copositive. To see this directly, maximize $\sum_{\{i,j\}\in E(C_5)}x_ix_j$ on $x\ge0$, $\sum x_i=1$, choosing a maximizer with smallest support. Two positive nonadjacent coordinates can be merged into the one with larger neighbor sum without reducing the objective, so the support is a clique. The cycle is triangle free, and the maximum over a clique is at most $1/4$. Hence $H_5[x]\ge0$.

Embed $H_5$ on the ordered coordinates $(1,2,3,4,n)$, zero elsewhere, and let $W_0=D^{-1}H_5D^{-1}$ (with this embedding understood). Then $W_0\in\COP_n$, $\norm{W_0}_F\le5$, and
\[
 \ip{W_0}{DFD}=\ip{H_5}{I_5+(3/5)\operatorname{Adj}(C_5)}=5-6=-1.
\]
The anchor changes this pairing by at most $\epsilon\norm{W_0}_F\norm C_F\le1/20$.
\end{proof}

\begin{lemma}[One strict separator for the whole perturbation ball]\label{lem:separator}
Define
\begin{equation}\label{eq:separator}
 W_N=S_N^{-1}W_0S_N^{-1}+\frac1{20nN^2}I.
\end{equation}
Every $A$ in \eqref{eq:ball} has all the positivity, norm, separation, and distance properties in Theorem~\ref{thm:main}.
\end{lemma}
\begin{proof}
Positive diagonal congruence preserves ordinary complete positivity and copositivity. The first term in \eqref{eq:separator} is copositive with norm at most five and pairing with $B_N$ at most $-19/20$. Since $\tr B_N\le3nN^2$, the identity addition increases that pairing by at most $3/20$. Thus
\[
 \ip{W_N}{B_N}\le-4/5,\quad
 \ip{W_N}{A_N}\le-1/(5nN^2),\quad \norm{W_N}_F\le6.
\]
The positive identity coefficient is the stated cone margin.

We have $\norm{A_N}_F\le3/4$, $\lambda_{\min}(A_N)>1/(160nN^2)$, and every entry of $A_N$ is at least $\epsilon/(4nN^2)=4\rho$. These facts survive the perturbation and give $\norm A_F\le1$. The pairing increases by at most $6\rho$, which leaves it at most $-1/(6nN^2)$. Finally, for any $X\in\CP_n$,
\[
 \norm{A-X}_F\ge\frac{-\ip{W_N}A}{\norm{W_N}_F}\ge\frac1{36nN^2}.
\]
\end{proof}

\section{All improving choices remain in the canonical face}
For an $A$ in the ball, multiply the objective by the positive factor $4nN^2$ and write
\[
 \widetilde B=4nN^2A=B_N+Z,\qquad \norm Z_F\le\epsilon/4.
\]
This scaling changes no improving direction.

\begin{lemma}[Arithmetic confinement]\label{lem:confinement}
At every vertex covered by Lemma~\ref{lem:heights}, every improving extreme tangent ray of the full $\Ry_n$ lies in the tangent cone of $\mathcal F$. Every such improving ray is bounded and has an adjacent endpoint in $\mathcal F$.
\end{lemma}
\begin{proof}
Normalize a ray $R$ as in \eqref{eq:rayheight}. Define
\[
 h(R)=\ip{B_0}{R_{[r],[r]}}.
\]
Each top-interval vector is active, so each of its quadratic values on $R$ is nonnegative. Lemma~\ref{lem:data} gives $h(R)\ge0$, with equality exactly when the top block of $R$ vanishes. If $h(R)>0$, integrality and the common denominator yield $h(R)\ge1/(500n^3)$. Hence
\begin{align*}
 \ip{\widetilde B}R
 &\ge\frac{N^2}{500n^3}-(4N+2)E-\frac{nE\epsilon}{4}\\
 &\ge\frac{N^2}{500n^3}-7NE
 \ge\frac{N^2}{1000n^3}>0,
\end{align*}
where the last step uses $N=20000n^3E^2$ and $E\ge1$. Thus every improving ray has top block zero.

The hyperplane $h=0$ exposes the tangent face corresponding to $\mathcal F$. Its rays are tangent rays of the epigraph \eqref{eq:face}. The recession directions there have nonnegative cross-column and bottom entry by Lemma~\ref{lem:recession}. Since $\widetilde B$ has strictly positive entries, none can be improving. Every improving extreme ray is therefore bounded and reaches a neighboring vertex in $\mathcal F$.
\end{proof}

\begin{lemma}[No vertex exit and no identity-corrected ray exit]\label{lem:noexit}
At every vertex covered by Lemma~\ref{lem:heights},
\[
 \ip{\widetilde B}P\ge N^2/4>0.
\]
For every improving extreme ray there, no real $\tau$ makes $R+\tau I$ a copositive separator of $A$.
\end{lemma}
\begin{proof}
Put $M=m+1$ and $\alpha=\ip{B_0}{T_r}\ge1$. Since $c>0$ and $a_0>0$,
\begin{align*}
 \ip{\widetilde B}P
 &\ge N^2-12NM-\frac\epsilon4\,10n^5M^2\\
 &=N^2-12NM-\frac{n^2M^2}{40}\ge N^2/4.
\end{align*}
The parameter choice implies $N\ge24M$ and $N\ge nM$.

For the second claim use the primitive integral generator, which is harmless under positive rescaling. By confinement,
\[
 R=\begin{pmatrix}0&w\\w\trans&s\end{pmatrix},\qquad |w_i|,|s|\le E.
\]
Strict positivity of the objective implies that an improving $R$ has $s<0$ or some $w_i<0$. Copositivity of $R+\tau I$ requires $\tau\ge0$. If $s<0$, it requires $\tau\ge-s\ge1$. Otherwise, for a negative integer $w_i$, its $(i,n)$ principal block requires
\[
 \tau(s+\tau)\ge w_i^2\ge1.
\]
If $\tau<1$, then $s+\tau\le E+1\le2E$, so in all cases $\tau\ge1/(2E)$. Also
\[
 \ip{\widetilde B}R\ge-7NE,\qquad \tr\widetilde B\ge N^2.
\]
It follows that
\[
 \ip{\widetilde B}{R+\tau I}
 \ge-7NE+\frac{N^2}{2E}>0,
\]
since $N>14E^2$. Therefore a copositive correction cannot separate. In particular the uncorrected improving ray is not copositive.
\end{proof}

\begin{proof}[Proof of Theorem~\ref{thm:main}]
Lemma~\ref{lem:separator} proves the target and certificate assertions. Starting at $T_n\in\mathcal F$, suppose fewer than $m$ pivots have occurred. Every previous improving step remains in $\mathcal F$ by confinement, so the movement theorem and Lemma~\ref{lem:heights} apply to the current vertex. The objective cannot be in its minimizing normal cone, because that cone is generated by nonnegative integer rank-one matrices and is contained in $\CP_n$, whereas $A\notin\CP_n$. Since the local tangent cone is polyhedral and pointed, it has an improving extreme ray. Every possible such ray is bounded and remains in $\mathcal F$, by Lemma~\ref{lem:confinement}. Neither allowed exit is possible, by Lemma~\ref{lem:noexit}. Induction therefore forces at least $m$ pivots, independently of all choices.

For fixed $n$, \eqref{eq:params} gives
\[
 N=K_n(m+1)^{4(d-1)},\qquad
 K_n=20000n^3\bigl((d-1)!\bigr)^2(8Y^2n^4)^{2(d-1)}.
\]
Every center entry has denominator dividing $2000n^4N^2$ and magnitude at most one, so $L_m=O_n(\log N)$. The distance lower bound bounds the reciprocal-distance term by $O_n(\log N)$ as well. Conversely the positive bottom entry is a fixed nonzero rational multiple of $N^{-2}$, whose reduced denominator has bit length $\Omega_n(\log N)$. Thus $H_m=\Theta_n(\log(m+1))$. The separator entries have logarithmic height by \eqref{eq:separator}. All its construction and verification identities require only polynomially many fixed-dimensional rational operations on these bit lengths.
\end{proof}

\section{Interpretation, prior work, and reproducibility}
The result is an ordinary-cone statement in every dimension where positive definite, strictly positive nonmembers are possible. There is no rectangular embedding of a constrained column cone, no noncanonical shear start, no assumed shadow trajectory, and no uniqueness premise. The fixed positive-definite-block face contains many possible trajectories; the bounded-movement theorem controls them all.

The two mechanisms are separate but compatible. Rank-one exposing normals and an integral objective gap confine all improving directions to a face. The geometry of a uniformly positive quadratic epigraph bounds the length of every edge in the remaining coordinates. Scaling the fixed nonmember produces an exponentially long unavoidable residence in that face. The same nonmember still has a directly available Horn-based separator. Thus local traversal complexity is not certificate-size complexity.

The simplex and neighboring-perfect-form framework is due to Dutour Sikiri\'c, Sch\"urmann, and Vallentin \cite{DSV}. Their Conjecture 3.2 concerns termination for a suitable pivot rule, not a polynomial pivot bound. The present theorem neither refutes finite termination nor claims a complexity lower bound for all separation algorithms. Dannenberg--Sch\"urmann \cite{DS} study the neighborhood graph and locally finite polyhedral structure; the canonical root-lattice form and that geometry are antecedents. Oertel--Sch\"urmann \cite{OS} generalize perfect-matrix theory and qualitative Interior Ryshkov results. Dirichlet approximation, determinant bounds, the Horn separator, and positive diagonal congruence are classical ingredients. The specific asserted contribution is the bounded-coordinate-movement lemma and its all-rule, canonical-start, ordinary-CP consequence.

The exact programs accompanying this paper verify the explicit target/separator arithmetic, the top-interval decomposition, all finite universal parameter inequalities used above, the canonical tangent rays, and sampled branching prefixes of the fixed-block face. A separate verifier imports no generator code and rejects deliberately corrupted data. Tests are supplements to the universal proof; no long trajectory is simulated. A proof over every pivot choice cannot be replaced by a finite sampling of choices. No independent human review, formal proof-assistant certification, exhaustive historical novelty certificate, or journal-percentile claim is made.

\begin{thebibliography}{9}
\bibitem{DSV} M. Dutour Sikiri\'c, A. Sch\"urmann, and F. Vallentin, A simplex algorithm for rational cp-factorization, \emph{Mathematical Programming} \textbf{187} (2021), 25--45. DOI: 10.1007/s10107-020-01467-4. arXiv:1807.01382v4.
\bibitem{DS} V. Dannenberg and A. Sch\"urmann, Perfect copositive matrices, \emph{Communications in Mathematics} \textbf{31} (2023), 137--156. DOI: 10.46298/cm.11141. arXiv:2303.17310v4.
\bibitem{OS} A. Oertel and A. Sch\"urmann, Generalized Perfect Matrices, arXiv:2602.05841v2 (2026).
\end{thebibliography}
\end{document}
